\section{Discussion}
\label{sec:discussion}

\subsection{Evidence helps globally but can hurt locally}
The practical headline is not that gold evidence fails.
It succeeds on average, and it rescues far more claim-only errors than it creates.
The scientific headline is that those averages are the wrong unit for reliability engineering.
A 1.15--1.51\% regression rate is small in a 10{,}000-claim table and large if each regression is a previously correct decision that evidence reversed.
Evaluation reports that publish only condition-wise accuracy will miss that reversal.

\subsection{What adjudication can and cannot say about regressions}
For roughly a third (Grok) to nearly half (Claude) of the 226 regressions, a blinded panel reaches the FEVER label from the evidence sentences alone and keeps it as titles and structured evidence are added.
On those items the designated evidence appears adequate to an independent reader while GPT moved away from the correct label after receiving it.
That pattern is compatible with an evidence-utilization failure and connects to work showing that models can ignore, mishandle, or be distracted by supplied context \citep{longpre-etal-2021-entity,shi2023distracted,wan-etal-2024-evidence}.
It does not establish that explanation: the original runs stored labels only, and the adjudicator is a different model reading the evidence in a different prompt.
A further fifth to quarter of regressions end in a decisive verdict opposite to FEVER, and a sixth (Claude) to nearly a third (Grok) end nondecisive.
In neither case does a decisive GPT answer after evidence show that GPT failed to use evidence supporting the FEVER answer.
These observations therefore do not establish the prevalence or cause of GPT evidence-utilization failures.

\subsection{Nondecisiveness and disagreement with FEVER}
Final nondecisiveness is common in both families, which is expected given prior evidence that sufficiency and ambiguity are first-class problems in fact verification \citep{atanasova-etal-2022-fact,glockner-etal-2024-ambifc}.
We do not interpret nondecisive outcomes or decisive judge--FEVER disagreements as FEVER annotation errors; inspected rationales show substantive conflicts (for example, a date or scope mismatch between claim and evidence), but judge rationales are observations, not labels.
Shared regressions, which both GPT models got wrong after evidence, are more often nondecisive under both families, although the Grok difference is not significant after correction.

\subsection{Disclosure and representation}
Titles and structured evidence change a minority of verdicts and never reverse a decisive consensus.
Stage~C adds alternative sentence text for some claims, so its effect mixes added information with changed representation.
Separating the two would require a structure-only arm that shows the chosen evidence set without alternatives; we did not run one.
We therefore make no claim about the size or portability of a pure representation effect.

\subsection{Judge-family dependence as measurement uncertainty}
The Claude-versus-Grok difference is a result, not a defect to be explained away.
LLM-as-judge work already shows that evaluators have family-specific biases and are not drop-in substitutes \citep{zheng2023judging,liu-etal-2023-g,chen-etal-2024-humans,huang-etal-2025-empirical}.
Here the same packets, rubric, and taxonomy produce a 13--16~pp difference in nondecisiveness at every stage and disagreement on about one regression in four.
Category shares are therefore properties of an adjudicating family as much as of the claims; two instruments do not define an uncertainty interval for a latent prevalence, and we report both rather than a single estimate.

\subsection{Implications for evidence-grounded LLM evaluation}
Four implications follow.
First, paired no-evidence versus gold-evidence evaluation should be standard whenever a benchmark supplies designated evidence; otherwise regressions are invisible.
Second, automated diagnostic categories need a second judge family before their shares are treated as stable, and decisiveness should be reported separately from agreement with the benchmark label.
Third, residual nondecisiveness should remain a first-class output of adjudication rather than being forced into Supported/Refuted.
Fourth, adjudication inputs deserve the same verification as model inputs: a field-extraction error silently corrupted the structured stage of our original adjudication.
% COAUTHOR REVIEW: the historical draft named 58 claim-level disagreements; the corrected count is 54 (accepted) or 58 (first-valid).
Human review of the 54--58 claim-level Grok--Claude category disagreements would be the natural next measurement, not another automated family.
